\section{Output-Length Sensitivity}
\label{app:length-diagnostic}

Pairwise evaluators may respond to output length as well as substantive
quality. We therefore test whether response length is associated with pairwise
performance. For each run $r$, task $t$, and regime $m$, we calculate the
Spearman correlation between an output's word count and its mean pairwise win
rate across eligible leave-family-out judges. Correlations are calculated
within each comparison set before being averaged across the ten runs, thereby
avoiding comparisons between tasks with different legitimate length demands.

Figure~\ref{fig:output-length-diagnostic} shows a positive association in most
settings. Averaged across tasks, the mean within-run correlation is $0.52$ in
automation and $0.27$ in augmentation. As an intuitive cross-check, the longer
output wins 72.5\% of decisive model-pair majorities in automation and 60.0\%
in augmentation. The pattern is not universal: the augmentation association is
weak for travel planning and tutoring and negative for menu planning. These
estimates are descriptive because model pairs share outputs and judges and are
therefore not independent observations.

\begin{figure}[H]
    \centering
    \includegraphics[width=\linewidth]{paper_figures/figureA_output_length_diagnostic.pdf}
    \caption{
    Output-length sensitivity across ten runs. Panel A reports the mean
    within-run Spearman correlation between output word count and mean pairwise
    win rate; error bars show $\pm 1.96$ standard errors across runs. Panel B
    reports the share of decisive model-pair majorities won by the longer
    output. The dashed line marks 50\%. Positive associations can reflect
    useful completeness, evaluator preference for elaboration, or both.
    }
    \label{fig:output-length-diagnostic}
\end{figure}

To assess sensitivity, we regress win rate on log word count separately within
each run--task--regime comparison set and re-rank the residualized outcomes.
This adjustment is deliberately conservative: length is partly an outcome of
model behavior and may mediate genuine completeness, so removing it can also
remove real quality. Nevertheless, the exercise shows that exact automation
rankings are length-sensitive. GPT-O4-Mini ranks first after adjustment,
followed by Claude-Opus-4.8 and GPT-5-Mini; Claude-Sonnet-4.6 falls markedly
because its automation outputs are substantially longer on average. In
augmentation, GPT-5-Mini remains the highest-ranked assistant model, while the
unaided worker becomes first overall. Most importantly for the paper's central
claim, the model-level correlation between automation and augmentation ranks
decreases from $\rho=0.48$ to $\rho=0.35$ after adjustment. Thus output length
affects the precise ordering, but it does not explain away the divergence
between direct-solving and assistance performance.
